% Section 4: measurements of the certified head on the engine's states (setup, recomputation,
% fallback and agreement, the engine check, head-only runtime). The file keeps its old name
% so that the editorial tools that list the main-text files keep working.
% Claims and status: pins from SETUP.md and the evidence READMEs; capture method and
% alignment from committed head-geometry evidence; held-out replay, kernel replay, engine
% check, head-only microbenchmark and primitive costs from committed evidence. No served run
% of the certified head is reported (pending bench-cert). Figure 4 (Section 5's) is placed in
% this file so that it floats onto the page where Section 5 begins.
\section{Measurements on the Engine's States}\label{sec:measure}
This section measures the certified head on states recorded from the running engine. Every number comes from one of the five sources of Table~\ref{tab:sources}; numbers labelled baseline are the stock head's.

\begin{figure*}[t]
\centering
% Figure fig:candidates. (a) reads evidence/head_geometry/selfevidence_plain4b_ccdf.csv.
% (b): the curves are derived, 1-(1-p)^B with the per-row interval-form shares of the
% held-out rows in evidence/head_geometry/rstock_plain4b.json (rules/gamma_1.19e-4/bucket/
% positions_needing_stock_kernel = 0.0034970857618651123 and rules/tensor_core_model/
% bucket/positions_needing_stock_kernel = 0.013988343047460449; these exist only in JSON,
% so they are copied here at full precision, and the legend's p and the shares at B = 128
% are computed from them). The markers are measured: the GPU kernel's share of batches of M
% consecutive rows of the first 60,000 capture rows that need a fallback
% (evidence/certified_head/head_path_time.csv, batch_fallback_rate and
% hopper_batch_fallback_rate).
% (c) reads evidence/certified_head/head_path_time.csv: the W8A16 path's time per call over
% the stock head and argmax (stock_us) on a decided batch (w8a16_us), and the expected time
% with the whole-batch fallback under each model (w8a16_expected_us,
% hopper_w8a16_expected_us) and with the column fallback under the conservative model
% (columns_expected_us). A microbenchmark of the head alone, probes on.
\pgfmathsetmacro{\phop}{0.0034970857618651123}
\pgfmathsetmacro{\pcons}{0.013988343047460449}
\pgfmathsetmacro{\fbhop}{100*(1-(1-\phop)^128)}
\pgfmathsetmacro{\fbcons}{100*(1-(1-\pcons)^128)}
\begin{tikzpicture}
\begin{groupplot}[group style={group size=3 by 1,horizontal sep=1.25cm},
  width=5.9cm,height=3.55cm,grid=major]
\nextgroupplot[xmode=log,ymode=log,xmin=1,xmax=248320,ymin=1e-4,ymax=1.5,
  xlabel={Candidate rows $k$},ylabel={Positions needing $\ge k$ rows},
  title={\paneltitle{a}{Candidates (measured)}},unbounded coords=discard,
  xtick={1,10,100,1000,10000,100000,248320},
  xticklabels={1,10,$10^2$,$10^3$,$10^4$,$10^5$,$V$},
  ytick={1,0.1,0.01,0.001,0.0001},yticklabels={100\%,10\%,1\%,0.1\%,0.01\%},
  minor tick num=0,xminorticks=false,yminorticks=false,
  legend style={at={(0.97,0.03)},anchor=south east},
  legend image post style={xscale=0.5}]
\addplot[series int8 row] table[x=k,y=plain_int8_row_greedy,col sep=comma]{../evidence/head_geometry/selfevidence_plain4b_ccdf.csv};\addlegendentry{int8 per row}
\addplot[series int8 g128] table[x=k,y=plain_int8_g128_greedy,col sep=comma]{../evidence/head_geometry/selfevidence_plain4b_ccdf.csv};\addlegendentry{int8 128 cols}
\addplot[vpblue,semithick,only marks,mark=o,mark size=1.4pt] table[x=k,y=plain_int8_g128_gumbel_t1.0,col sep=comma]{../evidence/head_geometry/selfevidence_plain4b_ccdf.csv};\addlegendentry{same, $T{=}1$}
\addplot[series fp8] table[x=k,y=plain_fp8_row_greedy,col sep=comma]{../evidence/head_geometry/selfevidence_plain4b_ccdf.csv};\addlegendentry{FP8 per row}
\addplot[series int4] table[x=k,y=plain_int4_g128_greedy,col sep=comma]{../evidence/head_geometry/selfevidence_plain4b_ccdf.csv};\addlegendentry{int4 128 cols}
\nextgroupplot[xmode=log,log basis x=2,xmin=1,xmax=256,ymin=0,ymax=1,
  xtick={1,4,16,64,256},xticklabels={1,4,16,64,256},
  xlabel={Rows in the batch $B$},ylabel={Batches with a fallback},
  ytick={0,0.2,0.4,0.6,0.8,1},yticklabels={0\%,20\%,40\%,60\%,80\%,100\%},
  title={\paneltitle{b}{Whole-batch fallback}},
  legend style={at={(0.03,0.97)},anchor=north west}]
\addplot[series conservative,domain=1:256,samples=120]{1-(1-\pcons)^x};\addlegendentry{cons., derived}
\addplot[series hopper,domain=1:256,samples=120]{1-(1-\phop)^x};\addlegendentry{Hopper, derived}
\addplot[only marks,mark=square*,mark size=1.5pt,vpgreen] table[x=M,y=batch_fallback_rate,col sep=comma]{../evidence/certified_head/head_path_time.csv};\addlegendentry{cons., kernel}
\addplot[only marks,mark=*,mark size=1.6pt,ink] table[x=M,y=hopper_batch_fallback_rate,col sep=comma]{../evidence/certified_head/head_path_time.csv};\addlegendentry{Hopper, kernel}
\draw[muted,densely dotted] (axis cs:128,0) -- (axis cs:128,1);
\node[annot,anchor=south east] at (axis cs:124,{\fbcons/100}) {\vppct{0}{\fbcons}};
\node[annot,anchor=south east] at (axis cs:122,{\fbhop/100+0.01}) {\vppct{0}{\fbhop}};
\nextgroupplot[xmode=log,log basis x=2,xmin=1,xmax=256,ymin=0.5,ymax=2.5,
  xtick={1,4,16,64,256},xticklabels={1,4,16,64,256},
  ytick={0.5,1,1.5,2,2.5},yticklabels={0.5,1,1.5,2,2.5},
  xlabel={Rows in the batch $M$},ylabel={Time over the stock head},
  title={\paneltitle{c}{Head-path time (measured)}},
  legend style={at={(0.03,0.97)},anchor=north west},
  every axis plot/.append style={mark size=1.4pt}]
\draw[muted,densely dotted] (axis cs:1,1) -- (axis cs:256,1);
\addplot[series int8 row,mark=*] table[x=M,y expr={\thisrow{w8a16_us}/\thisrow{stock_us}},col sep=comma]{../evidence/certified_head/head_path_time.csv};\addlegendentry{decided batch}
\addplot[series conservative,mark=square*,mark options={solid}] table[x=M,y expr={\thisrow{w8a16_expected_us}/\thisrow{stock_us}},col sep=comma]{../evidence/certified_head/head_path_time.csv};\addlegendentry{whole batch, cons.}
\addplot[series hopper,mark=triangle*] table[x=M,y expr={\thisrow{hopper_w8a16_expected_us}/\thisrow{stock_us}},col sep=comma]{../evidence/certified_head/head_path_time.csv};\addlegendentry{whole batch, Hopper}
\addplot[vpgreen,thick,dashdotted,mark=diamond*,mark options={solid}] table[x=M,y expr={\thisrow{columns_expected_us}/\thisrow{stock_us}},col sep=comma]{../evidence/certified_head/head_path_time.csv};\addlegendentry{columns, cons.}
\end{groupplot}
\end{tikzpicture}

\caption{What the certified head costs. \textbf{(a)}~Held-out replay: share of positions that need at least $k$ candidates, for four copies of the head (greedy) and for the 128-column int8 copy in a seeded race at $T=1$ \ev{ccdf}. \textbf{(b)}~Share of head calls of $M$ positions that fall back to the stock head: lines if positions fell back independently at the held-out shares of Table~\ref{tab:fallback} (derived) \ev{rstock}; markers measured in the kernel replay \ev{kernel-path}. \textbf{(c)}~Time per call over the stock head's, head alone: measured with every position certified (median of 30 replays; p10--p90 whiskers hidden by the markers below 256), and expected with each fallback under the \modelcons, from measured times and the rates of (b) \ev{kernel-path}; under the \modelhopper{} the expected times lie between these lines and the measured one. Shaded: above 16 positions the calls are consecutive replay positions, not engine batches.}
\label{fig:candidates}\label{fig:geometry-candidates}
\end{figure*}

\subsection{Setup}\label{sec:setup}\label{sec:transport-capture}
\paragraph{System} One NVIDIA GH200 with SGLang~\cite{sglangpaper} at commit \sourcepin{}, PyTorch 2.13 for CUDA 13.0 and the FlashInfer attention backend (Appendix~\ref{app:pins} lists every version). The target is Qwen3.5-4B at revision \modelpin~\cite{qwencard}, whose layers mix linear attention (Gated DeltaNet~\cite{gdn}) with full attention. The drafters are MTP-4B, the model's native MTP layer (four draft steps, top-$k$ 1), and DFlash-4B, the public block drafter for this target~\cite{dflashpaper,dflash4bcard}, which drafts blocks of 16 tokens and projects them through the target's head.

\paragraph{Capture} A capture-only engine patch records each position's head input as BF16 bit patterns, with the token the engine chose; under speculation it also records the drafter's head inputs and pairs each drafted token with the target head input that verifies it. The capture servers ran without CUDA graphs or the overlap scheduler, with at most 16 running requests in plain decoding (Appendix~\ref{app:transport-capture}). The 320 \emph{capture prompts} come from eight public datasets in four domains and were decoded greedily for 384 tokens; they were split by prompt into analysis and held-out halves before anything was fitted. The capture records what the head read: in plain decoding the FP64 argmax of the captured head input equals the engine's token at 99.51\% of positions, and every disagreement lies within one BF16 spacing of the FP64 winner \ev{alignment}.

\begin{table}[t]
\caption{Where the numbers come from. A replay re-evaluates recorded head inputs offline; the held-out replay draws whole decode steps (750) at random from the held-out half until it has 6,000 positions.}\label{tab:sources}
\vskip 0.05in
\centering\scriptsize
\begin{tabularx}{\columnwidth}{@{}>{\raggedright\arraybackslash}p{1.55cm}Y>{\raggedright\arraybackslash}p{1.6cm}@{}}
\toprule
Source & What & Size\\\midrule
Held-out replay & FP64 replay of held-out plain-decoding positions & 6,005 positions\\\addlinespace[2pt]
Kernel replay & the GPU kernel on the first captured positions, at their batch shapes & 60,000 positions\\\addlinespace[2pt]
Engine check & the kernel inside SGLang, \modelcons & 64 prompts\\\addlinespace[2pt]
Microbenchmark & the head alone, CUDA graphs, warm L2 cache & $M=1$ to 256\\\addlinespace[2pt]
Speculative captures & aligned draft and target head inputs; verification positions & 40,000 pairs per drafter\\
\bottomrule
\end{tabularx}
\end{table}

\subsection{How Many Entries Are Recomputed}\label{sec:geometry-selfevidence}\label{sec:candidates}
In the kernel replay the certified head recomputes 1.84 entries per decision on average under the \modelcons{} and 1.78 under the \modelhopper; the median is 1 and the 99th percentile 9 \ev{kernel-replay}. The held-out replay compares copies of the head on the simpler \Rreal{} decision, whose screen needs no margin \ev{selfevidence,selfevidence-csv}. There, no error bound was violated and the exact winner was always a candidate. The int8 copy with one scale per entry needs 1.55 candidates on average, and at 75.4\% of positions a single one. With one scale per 128 columns the mean falls to 1.34, at 0.52 of the BF16 head's bytes (Figure~\ref{fig:candidates}a; Table~\ref{tab:candidates} gives percentiles). FP8 E4M3 needs 4.9 on average, and int4 decides no position at all. Int8 suffices because of the margins, not because its bound is tight (Section~\ref{sec:envelope}): its intervals are about two logits wide, while the median gap between the two largest logits is 7.3 logits \ev{stats}, so tightening the bound helps little (Appendix~\ref{app:candidates}). On the verification positions of the speculative captures the per-entry copy needs 1.57 candidates for DFlash-4B and 1.51 for MTP-4B \ev{selfevidence-dflash,selfevidence-mtp}.

\begin{figure*}[t]
\centering
% Figure fig:rho. (a) is drawn to scale from medians: ||h^d|| = 50.596 and ||h^t|| = 159.2
% (evidence/head_geometry/transport_dflash4b.json, headline hdnorm2/htnorm2 medians),
% cos(h^d, h^t) = 0.408 (evidence/head_geometry/stats_dflash4b.json, pairs/cos_hd_ht/all
% q50), and the median threshold delta* = 0.00854 (evidence/precision/head_constants.json,
% tiles/contiguous/euclid_int8_row_threshold/p50). (b) reads
% evidence/head_geometry/rho_dflash4b_quantiles.csv; the MTP-4B markers are the p10,
% median and p90 of evidence/head_geometry/transport_mtp4b.json,
% headline/dnorm2_over_htnorm2 (0.7666, 0.9541, 1.1756).
\begin{tikzpicture}[x=0.0335cm,y=0.0335cm,font=\footnotesize]
\pgfmathsetmacro{\Lt}{159.2}
\pgfmathsetmacro{\Ld}{50.596}
\pgfmathsetmacro{\Cth}{0.408}
\pgfmathsetmacro{\Sth}{sqrt(1-\Cth*\Cth)}
\coordinate (O) at (0,0);
\coordinate (T) at (\Lt,0);
\coordinate (D) at ({\Ld*\Cth},{\Ld*\Sth});
\coordinate (F) at ({\Lt*\Cth*\Cth},{\Lt*\Cth*\Sth});
\path[use as bounding box] (-40,-40) rectangle (176,104);
\node[anchor=north,font=\footnotesize] at (68,104) {(a) Median DFlash-4B head inputs, to scale};
% the line through h^d and the foot of the perpendicular from h^t (best rescaling of h^d)
\draw[muted,densely dotted] (O) -- ($(O)!1.18!(F)$);
\draw[muted,densely dotted] (T) -- (F);
\node[note,anchor=west,align=left] at ($(F)+(3,4)$) {closest rescaling of $h^d$:\\$\rho\approx0.91$};
% vectors
\draw[flow,very thick,shorten >=2.2pt] (O) -- (T);
\node[below,text=ink] at (80,-2) {$h^t$, $\norm{h^t}_2\approx159$};
\draw[flow,very thick,teal] (O) -- (D);
\node[text=teal,anchor=east,align=right] at ($(D)+(-3,-10)$) {$h^d$, $\norm{h^d}_2$\\$=\sqrt{2560}$};
\draw[flow,thick,brick,dashed,shorten >=2.2pt] (D) -- (T);
\node[text=brick,anchor=south west] at (78,22) {$\Delta$: $\norm\Delta_2\approx0.92\norm{h^t}_2$};
\draw[muted] ($(O)+(13,0)$) arc[start angle=0,end angle=65.9,radius=13];
\node[note,anchor=west] at (14,5) {$\cos=0.41$};
% the ball inside which transport would beat int8: radius 0.00854*159.2 = 1.36 units, to scale
\fill[ochre] (T) circle[radius={0.00854*\Lt}];
\draw[ochre,thin] (T) circle[radius=2.6pt];
\draw[ochre,thin] (T) -- (152,-14);
\node[note,text=ochre,anchor=north,align=center] at (150,-14) {transport beats int8 only\\inside the dot in this circle:\\radius $0.0085\norm{h^t}_2$};
\end{tikzpicture}\hfill
\begin{tikzpicture}
\begin{semilogxaxis}[width=8.9cm,height=4.7cm,xmin=2e-3,xmax=1.6,ymin=0,ymax=1,
  xlabel={Drift ratio $\rho$ of a pair, or per-row threshold $\delta^\ast_i$},ylabel={Cumulative share},grid=major,
  title={(b) Measured drift against the thresholds},xminorticks=false,
  legend style={at={(0.5,-0.36)},anchor=north,legend columns=3,
    /tikz/every even column/.append style={column sep=5pt}}]
\addplot[thick,teal,dashed] table[x=threshold_int8_row,y=q,col sep=comma]{../evidence/head_geometry/rho_dflash4b_quantiles.csv};\addlegendentry{$\delta^\ast_i$, int8 per row}
\addplot[thick,teal,densely dotted] table[x=threshold_int8_g128,y=q,col sep=comma]{../evidence/head_geometry/rho_dflash4b_quantiles.csv};\addlegendentry{$\delta^\ast_i$, int8 128 cols}
\addplot[thick,ochre,dashdotted] table[x=threshold_int4_row,y=q,col sep=comma]{../evidence/head_geometry/rho_dflash4b_quantiles.csv};\addlegendentry{$\delta^\ast_i$, int4 per row}
\addplot[very thick,ink] table[x=rho_all,y=q,col sep=comma]{../evidence/head_geometry/rho_dflash4b_quantiles.csv};\addlegendentry{$\rho$, DFlash-4B pairs}
\addplot[only marks,mark=triangle*,mark size=2.6pt,brick] coordinates {(0.7666151341559915,0.1) (0.954136519993249,0.5) (1.1756316575526062,0.9)};\addlegendentry{$\rho$, MTP-4B p10, 50, 90}
\draw[{Stealth[length=1.6mm]}-{Stealth[length=1.6mm]},muted] (axis cs:0.0095,0.5) -- (axis cs:0.84,0.5);
\node[font=\scriptsize,text=ink,fill=white,inner sep=1pt] at (axis cs:0.037,0.5) {medians $\approx100\times$ apart};
\end{semilogxaxis}
\end{tikzpicture}

\caption{Why transport fails. \textbf{(a)}~Median DFlash-4B head inputs, to scale (norms over the 4,020 transport pairs, cosine over 20,000 drift pairs). Transport's bound beats the int8 copy's only if $h^d$ lies in the dot at the tip of $h^t$, of radius $\delta^\ast\norm{h^t}_2$ for the median $\delta^\ast$ (Theorem~\ref{thm:crossover}); the best rescaling of $h^d$ (dotted) leaves $\rho=0.91$. \textbf{(b)}~Cumulative distributions of the measured drift $\rho$ over 40,000 aligned held-out pairs of each drafter, and of the threshold $\delta^\ast_i$ derived from the weights for each of the 248,320 vocabulary entries (64-entry tiles of contiguous token indices) \ev{head-constants,rho,rho-mtp,transport,stats-dflash}.}
\label{fig:rho}\label{fig:transport}
\end{figure*}

\subsection{Fallback and Agreement}\label{sec:rstock-share}
\begin{table}[t]
\caption{Share of positions the contract leaves to the stock head, by rule and by the error model assumed for the stock kernel. ``Steps'': share of the capture's decode steps (8.0 positions each) with such a position. The last row is a sensitivity check, not a model of the stock kernel. In full: Table~\ref{tab:rstock}.}\label{tab:fallback}
\vskip 0.05in
\centering\scriptsize\setlength{\tabcolsep}{3.5pt}
\begin{tabular*}{\columnwidth}{@{}l@{\extracolsep{\fill}}rrrr@{}}
\toprule
& \multicolumn{3}{c}{Held-out replay \ev{rstock}} & Kernel replay\\\cmidrule(lr){2-4}
Model ($\gamma$) & Gap & Interval & Steps & \ev{kernel-replay}\\\midrule
Conservative ($\gcons$) & 3.03\% & 1.40\% & 10.9\% & 1.46\%\\
Hopper ($\ghop$) & 2.18\% & 0.35\% & 2.8\% & 0.28\%\\
IEEE FP32 tree ($1.7\times10^{-6}$) & 1.95\% & 0\% & 0\% & --\\
\bottomrule
\end{tabular*}
\end{table}

Table~\ref{tab:fallback} gives how often each rule leaves a position to the stock head. The gap condition never falls below about 2\%, because it demands a margin of more than one BF16 spacing; the interval form leaves 0.35\% (\modelhopper) or 1.40\% (\modelcons) of the held-out positions. For every model, no position the interval form certified had a token different from the engine's at the capture's batch shape \ev{rstock}, and with the fallback the kernel replay matched the engine at all 60,000 positions \ev{kernel-replay}. This agreement cannot tell the models apart: even an IEEE FP32 summation tree, with a $\gamma$ about 360 times smaller than the conservative one, certifies every held-out position and gets every token right. It is consistent with each model, not a proof of any; the observed error of the stock FP32-output kernel, 58 times inside the \modelhopper, is the more direct test (Section~\ref{sec:limits}).

Fallback grows with the batch, because one undecided position is enough to send a whole batch to the stock head. If positions fell back independently with probability $p$, a call of $M$ positions would contain an undecided one with probability $1-(1-p)^M$. At the capture's 8.0 positions per step this matches the measured share of steps with a fallback, 2.8\% under the \modelhopper{} and 10.9\% under the \modelcons. At $M=128$ it predicts 36\% and 84\%, and the kernel replay measures 27\% and 77\% on calls of 128 consecutive positions (Figure~\ref{fig:candidates}b). A whole-batch fallback would therefore run the stock head for most large batches, so at serving batch sizes the certified head depends on the per-position fallback.

\paragraph{In the engine}\label{sec:engine} The engine check ran the kernel inside SGLang, untimed, on plain decoding, greedy verification and the drafters' projections, with the whole-batch fallback on every path but one plain-decoding run \ev{kernel-engine}. No certified position differed from the stock head on these runs, and with one request at a time the certified and stock servers' outputs were identical on all 64 prompts, for plain decoding and for MTP-4B. With up to 16 requests, 1.38\% of plain-decoding positions fell back (the kernel replay left 1.46\%, on other positions), and so did 1.74\% of MTP-4B verification positions (8 requests). DFlash-4B verification, with up to 128 positions per call, left 3.99\% of positions open, and the whole-batch fallback ran in 90\% of its verify head calls. We infer that there the per-position fallback, which reruns only an open position's candidates, would have to carry the result; it was exercised only in that plain-decoding run.

\subsection{Head-Only Runtime}\label{sec:kernel-runtime}
The certified path pays off only if
\begin{equation}
T_{\rm pass}+T_{\rm select}+T_{\rm rescore}+P_{\rm fb}(M)\,T_{\rm stock}(M)<T_{\rm chain}(M),\label{eq:headgate}
\end{equation}
where the left side adds the int8 pass, the selection of candidates, their recomputation and the stock head's time weighted by the share $P_{\rm fb}(M)$ of calls that fall back, and $T_{\rm chain}$ is the stock head itself: its GEMM, the FP32 copy and the argmax. As a baseline, the stock GEMM at $M=1$ runs at 95\% of the measured read bandwidth \ev{profile-head,profile-bw}, so a pass reading half the bytes needs at least about 170~\textmu s (derived).

In the microbenchmark a fully certified batch at $M=1$ takes 239.6~\textmu s, against 370.5~\textmu s for the stock head. Adding each fallback's measured cost at its measured rate gives an expected time of 0.67--0.70 of the stock head's up to 16 positions and 0.93 at 64, with the per-position fallback under the \modelcons{} \ev{kernel-path,kernel-rowwise}. From 128 positions every mode is slower (Figure~\ref{fig:candidates}c), because the int8 pass alone is slower than the stock head there; with the whole-batch fallback, 77\% of calls also pay for the stock head. The pass is not bound by memory bandwidth: at one position its directed-rounding epilogue keeps the streaming multiprocessors 79\% busy while it draws 62--63\% of DRAM bandwidth, so its time grows with $M$ faster than the stock GEMM's \ev{kernel-micro}. Since the head takes 10.3\% of a plain decode step at batch 1, an expected 0.67--0.70 of the stock head's time would shorten that step by at most about 3\% (derived, not measured). The effect on served decoding is not yet measured \pend{bench-cert}.
